% =====================================================================
%  FICHE 8 — THÈME 5 — NIVEAU FACILE — CORRIGÉ
% =====================================================================
\documentclass[11pt,a4paper]{article}
\usepackage[utf8]{inputenc}
\usepackage[T1]{fontenc}
\usepackage{lmodern}
\usepackage[french]{babel}
\usepackage{amsmath,amssymb,mathtools}
\usepackage{icomma}
\usepackage{siunitx}
\sisetup{output-decimal-marker={,},per-mode=power,group-separator={\,},group-minimum-digits=5}
\usepackage[version=4]{mhchem}
\usepackage{xcolor}
\usepackage{tikz}
\usepackage[most]{tcolorbox}
\usepackage{enumitem}
\usepackage{array}
\usepackage{fancyhdr}
\usepackage[hidelinks]{hyperref}
\usetikzlibrary{arrows.meta,positioning,calc}
\usepackage[a4paper,margin=2.1cm,headheight=26pt,headsep=10pt]{geometry}
\usepackage{titlesec}

\definecolor{primary}{RGB}{150,98,162}
\definecolor{primarydark}{RGB}{92,52,110}
\definecolor{excol}{RGB}{0,135,90}

\tcbset{boxbase/.style={enhanced,breakable,boxrule=0.9pt,arc=2.5pt,left=3mm,right=3mm,top=2.3mm,bottom=2.3mm,fonttitle=\bfseries,coltitle=white,toptitle=0.6mm,bottomtitle=0.6mm}}
\newtcolorbox[auto counter]{correction}[1][]{boxbase,colback=excol!4,colframe=excol,title={\textsc{Corrigé}~\thetcbcounter\ \ifx\relax#1\relax\relax\else\textmd{--- \itshape #1}\fi}}

\newcommand{\rep}[1]{\textcolor{primarydark}{\textbf{#1}}}

\pagestyle{fancy}
\fancyhf{}
\fancyhead[L]{\small\textcolor{primarydark}{\textbf{Fiche 8} — Thème 5 : Équilibre chimique et Le Chatelier}}
\fancyhead[R]{\small\textcolor{primarydark}{\textsc{Facile — corrigé}}}
\fancyfoot[C]{\small\thepage}
\renewcommand{\headrulewidth}{0.8pt}
\renewcommand{\headrule}{\hbox to\headwidth{\color{primary}\leaders\hrule height \headrulewidth\hfill}}
\setlength{\parindent}{0pt}\setlength{\parskip}{3pt}

\begin{document}

\begin{center}
\begin{tikzpicture}
\node[rectangle,rounded corners=7pt,fill=excol,text=white,minimum width=16cm,
      minimum height=1.1cm,align=center,inner sep=6pt]
  {{\large\bfseries Thème 5 — Équilibre chimique et Le Chatelier}\\[1pt]{\small Niveau \textsc{Facile} — corrigé détaillé}};
\end{tikzpicture}
\end{center}

\begin{correction}[écrire la constante d'équilibre]
Produits au numérateur, exposants = coefficients.
\begin{enumerate}[label=\alph*),leftmargin=2em,topsep=2pt,itemsep=3pt]
  \item $K = \dfrac{[\ce{NH3}]^2}{[\ce{N2}]\,[\ce{H2}]^3}$.
  \item $K = \dfrac{[\ce{PCl3}]\,[\ce{Cl2}]}{[\ce{PCl5}]}$.
  \item $K = \dfrac{[\ce{HI}]^2}{[\ce{H2}]\,[\ce{I2}]}$.
\end{enumerate}
\end{correction}

\begin{correction}[équilibres hétérogènes : que garde-t-on ?]
On retire les solides et liquides purs.
\begin{enumerate}[label=\alph*),leftmargin=2em,topsep=2pt,itemsep=3pt]
  \item $K = [\ce{CO2}]$ \; (les solides $\ce{CaCO3}$ et $\ce{CaO}$ disparaissent).
  \item $K = \dfrac{[\ce{C}]^2\,[\ce{D}]}{[\ce{A}]^2}$ \; (le solide $\ce{B}$ disparaît).
  \item $K = \dfrac{[\ce{CO2}]}{[\ce{CO}]}$ \; (les solides $\ce{FeO}$ et $\ce{Fe}$ disparaissent).
\end{enumerate}
\end{correction}

\begin{correction}[interpréter la valeur de $K$]
\begin{enumerate}[label=\alph*),leftmargin=2em,topsep=2pt,itemsep=2pt]
  \item $K=\num{1e5}\gg1$ : équilibre très déplacé vers les \rep{produits} (réaction quasi totale).
  \item $K=\num{2e-4}\ll1$ : équilibre très déplacé vers les \rep{réactifs} (réaction très limitée).
  \item $K=1{,}5\approx1$ : réactifs et produits \rep{coexistent} en proportions comparables.
\end{enumerate}
\end{correction}

\begin{correction}[une perturbation à la fois]
Le système s'oppose à la perturbation.
\begin{enumerate}[label=\alph*),leftmargin=2em,topsep=2pt,itemsep=2pt]
  \item Ajout de $\ce{N2}$ (réactif) : l'équilibre le consomme, déplacement vers la \rep{droite}.
  \item Retrait de $\ce{NH3}$ (produit) : le système en reforme, déplacement vers la \rep{droite}.
  \item $p\nearrow$ : vers le côté à \emph{moins} de gaz ; à gauche $4$ mol, à droite $2$ mol, donc vers la
        \rep{droite}.
  \item Catalyseur : \rep{aucun déplacement} (il ne fait qu'accélérer l'atteinte de l'équilibre).
\end{enumerate}
\end{correction}

\begin{correction}[la température change $K$]
\begin{enumerate}[label=\alph*),leftmargin=2em,topsep=2pt,itemsep=2pt]
  \item Exothermique : augmenter $T$ déplace vers le sens endothermique (les réactifs), donc $K$ \rep{diminue}.
  \item Endothermique : augmenter $T$ déplace vers les produits, donc $K$ \rep{augmente}.
  \item Seule la \rep{température} modifie la valeur de $K$ ; les concentrations initiales et le catalyseur
        ne la changent pas.
\end{enumerate}
\end{correction}

\end{document}
